\documentclass[12pt,a4paper]{article}
\usepackage[utf8]{inputenc}
\usepackage[T2A]{fontenc}
\usepackage[bulgarian]{babel}
\usepackage{geometry}
\usepackage{graphicx}
\usepackage{listings}
\usepackage{xcolor}
\usepackage{hyperref}
\usepackage{booktabs}
\usepackage{float}
\usepackage{tikz}
\usetikzlibrary{shapes,arrows,positioning}

\geometry{margin=2.5cm}

\lstset{
    basicstyle=\ttfamily\small,
    breaklines=true,
    frame=single,
    numbers=left,
    numberstyle=\tiny,
    keywordstyle=\color{blue},
    commentstyle=\color{green!60!black},
    stringstyle=\color{red}
}

\title{
    \vspace{2cm}
    \textbf{Курсов проект} \\
    \vspace{1cm}
    \Large{P2P Архитектури при Разпределено Динамично Балансиране на Паралелни Приложения}
}

\author{
    Изготвил: Георги Лазов \\
    Дисциплина: Разпределени ИТ Архитектури \\
    Специалност: Информационни системи
}

\date{Април 2026}

\begin{document}

\maketitle
\newpage

\tableofcontents
\newpage

\section{Въведение}

Паралелните и разпределени системи са ключови за решаването на изчислително интензивни задачи в съвременния свят. Един от основните предизвикателства при такива системи е ефективното разпределение на натоварването между процесорите/възлите, особено когато работното натоварване е нехомогенно или непредвидимо.

В централизираните архитектури (Master-Slave) един главен процес координира разпределението на задачите. Този подход е прост за имплементация, но създава ``тясно място'' (bottleneck) при мащабиране. Алтернативата е \textbf{Peer-to-Peer (P2P)} архитектура, където процесите комуникират директно помежду си без централен координатор.

Целта на този проект е да се изследва архитектурата на P2P паралелни приложения при разпределено динамично балансиране, като се анализират различни топологии на процесите - по-конкретно \textbf{верижна топология (chain)} и \textbf{пълен граф (all-to-all)}.

\section{Теоретична основа}

\subsection{Паралелни модели и SPMD}

При паралелното програмиране се различават два основни модела:

\begin{itemize}
    \item \textbf{SPMD (Single Program Multiple Data)} - всички процеси изпълняват една и съща програма, но върху различни данни. Това е най-разпространеният модел при MPI приложения.
    \item \textbf{MPMD (Multiple Program Multiple Data)} - различните процеси могат да изпълняват различни програми. Типичен пример са конвейрите за пакетна обработка.
\end{itemize}

В контекста на P2P динамичното балансиране работим основно с \textbf{SPMD модел}, където всички процеси са равнопоставени и могат да поемат произволна част от работата.

\subsection{Peer-to-Peer архитектури}

При P2P архитектурите няма централен координатор - всички възли (peers) са равнопоставени и комуникират директно помежду си.

\textbf{Предимства:}
\begin{itemize}
    \item Няма единична точка на отказ
    \item По-добра скалируемост
    \item Разпределено вземане на решения
\end{itemize}

\textbf{Недостатъци:}
\begin{itemize}
    \item По-сложна координация
    \item Възможно неоптимално разпределение
    \item Комуникационни разходи при пълен граф
\end{itemize}

\subsection{Топологии на процесите}

\subsubsection{Верижна топология (Chain/Ring)}

При верижната топология всеки процес комуникира само със съседите си (ляв и десен).

\textbf{Характеристики:}
\begin{itemize}
    \item Брой връзки: $n-1$ (за верига) или $n$ (за пръстен)
    \item Комуникационна сложност: $O(n)$ за глобална информация
    \item \textbf{Скалируем} - добавянето на нов процес изисква промяна само на съседните връзки
\end{itemize}

\subsubsection{Пълен граф (All-to-All)}

При пълния граф всеки процес може да комуникира директно с всеки друг.

\textbf{Характеристики:}
\begin{itemize}
    \item Брой връзки: $n(n-1)/2$
    \item Комуникационна сложност: $O(1)$ за директна комуникация
    \item \textbf{Не е скалируем} - комуникационните разходи растат квадратично
\end{itemize}

\subsection{Динамично балансиране на натоварването}

При динамичното балансиране работата се преразпределя по време на изпълнение. В P2P контекст използваме алгоритъма \textbf{Work Stealing}:

\begin{enumerate}
    \item Процес завършва задача
    \item Ако опашката му е празна, поисква работа от съсед
    \item Ако съседът има достатъчно работа, споделя половината
    \item Ако съседът също няма работа, продължава по веригата
\end{enumerate}

\section{Анализ на архитектурите}

\subsection{Централизирано vs. Разпределено балансиране}

\begin{table}[H]
\centering
\begin{tabular}{lcc}
\toprule
\textbf{Характеристика} & \textbf{Централизирано} & \textbf{Разпределено (P2P)} \\
\midrule
Скалируемост & Ограничена & По-добра \\
Сложност & Ниска & По-висока \\
Глобална информация & Да & Не (локална) \\
Отказоустойчивост & Ниска & По-висока \\
\bottomrule
\end{tabular}
\caption{Сравнение на подходите за балансиране}
\end{table}

\subsection{Сравнение на топологии}

\begin{table}[H]
\centering
\begin{tabular}{lcc}
\toprule
\textbf{Критерий} & \textbf{Верига} & \textbf{Пълен граф} \\
\midrule
Брой връзки & $O(n)$ & $O(n^2)$ \\
Латентност & $O(n)$ & $O(1)$ \\
Скалируемост & Висока & Ниска \\
Балансиране & Локално & Глобално \\
\bottomrule
\end{tabular}
\caption{Сравнение на топологиите}
\end{table}

\subsection{Хардуерни аспекти}

При облачни инфраструктури (напр. Amazon AWS) топологията на мрежата влияе пряко върху производителността:

\begin{itemize}
    \item \textbf{Placement Groups} - възлите в една група са с ниска латентност (<1ms)
    \item \textbf{Elastic Fabric Adapter} - специализиран интерфейс за HPC (до 100 Gbps)
    \item \textbf{NUMA архитектура} - трябва да се съобрази при многопроцесорни възли
\end{itemize}

\section{Проектиране}

\subsection{Функционално проектиране}

В рамките на проекта са реализирани три варианта:
\begin{enumerate}
    \item Статично циклично балансиране (baseline)
    \item Динамично P2P балансиране с верижна топология
    \item Динамично P2P балансиране с пълен граф
\end{enumerate}

\textbf{Проблем за тестване:} Рекурсивно изчисление на числата на Фибоначи - задача с нехомогенно натоварване.

\subsection{Архитектура на решението}

\textbf{Модел:} SPMD - всички процеси изпълняват един и същ код.

\textbf{Комуникационен протокол:}
\begin{itemize}
    \item \texttt{WORK\_REQUEST} - искане за работа
    \item \texttt{WORK\_RESPONSE} - отговор с работа
    \item \texttt{TERMINATE} - сигнал за край
\end{itemize}

\subsection{UML диаграми}

\subsubsection{Диаграма на последователността при верижна топология}

\begin{figure}[H]
\centering
\begin{tikzpicture}[node distance=2cm, auto]
    \tikzstyle{process} = [rectangle, draw, text width=1.5cm, text centered, minimum height=0.8cm]
    \tikzstyle{arrow} = [thick,->,>=stealth]

    \node[process] (P0) {P0};
    \node[process, right=of P0] (P1) {P1};
    \node[process, right=of P1] (P2) {P2};
    \node[process, right=of P2] (P3) {P3};

    \draw[arrow] (P1) -- node[above] {REQUEST} (P2);
    \draw[arrow, dashed] (P2) -- node[below] {RESPONSE} (P1);

    \draw[arrow] (P0) -- node[above, yshift=0.5cm] {REQUEST} (P1);
    \draw[arrow, dashed] (P1) -- node[below, yshift=-0.5cm] {NO\_WORK} (P0);
\end{tikzpicture}
\caption{Комуникация при верижна топология}
\end{figure}

\subsubsection{Диаграма на състоянията}

\begin{figure}[H]
\centering
\begin{tikzpicture}[node distance=2.5cm, auto]
    \tikzstyle{state} = [ellipse, draw, text width=2cm, text centered, minimum height=1cm]
    \tikzstyle{arrow} = [thick,->,>=stealth]

    \node[state] (init) {INIT};
    \node[state, right=of init] (working) {WORKING};
    \node[state, below=of working] (stealing) {STEALING};
    \node[state, right=of working] (term) {TERMINATED};

    \draw[arrow] (init) -- node[above] {start} (working);
    \draw[arrow] (working) edge[loop above] node[above] {task done} (working);
    \draw[arrow] (working) -- node[right] {empty queue} (stealing);
    \draw[arrow] (stealing) -- node[left] {got work} (working);
    \draw[arrow] (stealing) -- node[below] {no work} (term);
\end{tikzpicture}
\caption{Състояния на процес}
\end{figure}

\section{Имплементация}

\subsection{Технологичен избор}

\begin{itemize}
    \item \textbf{Език:} C
    \item \textbf{Паралелизъм:} MPI (Message Passing Interface)
    \item \textbf{Компилатор:} mpicc (OpenMPI)
\end{itemize}

\subsection{Програмен код}

Основните файлове са:
\begin{itemize}
    \item \texttt{static\_balancing.c} - статично циклично балансиране
    \item \texttt{p2p\_chain.c} - P2P с верижна топология
    \item \texttt{p2p\_full.c} - P2P с пълен граф
\end{itemize}

\section{Тестова среда}

\textbf{Хардуерна конфигурация:}
\begin{itemize}
    \item CPU: Apple M1 Pro (8 ядра)
    \item RAM: 16 GB
    \item L1 Cache: 192 KB (per P-core)
    \item L2 Cache: 12 MB (споделен)
\end{itemize}

\textbf{Параметри:}
\begin{itemize}
    \item Брой процеси (p): 1, 2, 4, 8
    \item Fibonacci стойност: 38
    \item Повторения: 3
\end{itemize}

\section{Резултати}

\subsection{Статично балансиране (Baseline)}

\begin{table}[H]
\centering
\begin{tabular}{cccccc}
\toprule
\textbf{p} & \textbf{$T_p$ (ms)} & \textbf{$S_p$} & \textbf{$E_p$} \\
\midrule
1 & 4498 & 1.00 & 1.00 \\
2 & 2876 & 1.56 & 0.78 \\
4 & 1798 & 2.50 & 0.63 \\
8 & 1423 & 3.16 & 0.40 \\
\bottomrule
\end{tabular}
\caption{Резултати при статично балансиране}
\end{table}

\subsection{P2P Верижна топология}

\begin{table}[H]
\centering
\begin{tabular}{cccccc}
\toprule
\textbf{p} & \textbf{$T_p$ (ms)} & \textbf{$S_p$} & \textbf{$E_p$} \\
\midrule
1 & 4498 & 1.00 & 1.00 \\
2 & 2398 & 1.88 & 0.94 \\
4 & 1298 & 3.47 & 0.87 \\
8 & 823 & 5.47 & 0.68 \\
\bottomrule
\end{tabular}
\caption{Резултати при P2P верижна топология}
\end{table}

\subsection{P2P Пълен граф}

\begin{table}[H]
\centering
\begin{tabular}{cccccc}
\toprule
\textbf{p} & \textbf{$T_p$ (ms)} & \textbf{$S_p$} & \textbf{$E_p$} \\
\midrule
1 & 4498 & 1.00 & 1.00 \\
2 & 2367 & 1.90 & 0.95 \\
4 & 1245 & 3.61 & 0.90 \\
8 & 912 & 4.93 & 0.62 \\
\bottomrule
\end{tabular}
\caption{Резултати при P2P пълен граф}
\end{table}

\section{Анализ на резултатите}

\subsection{Основни наблюдения}

\begin{enumerate}
    \item \textbf{Динамичното балансиране превъзхожда статичното} при 8 процеса:
    \begin{itemize}
        \item Статично: $S_p = 3.16$
        \item P2P Chain: $S_p = 5.47$ (\textbf{73\% подобрение})
        \item P2P Full: $S_p = 4.93$ (\textbf{56\% подобрение})
    \end{itemize}

    \item \textbf{Верижната топология показва по-добри резултати} при 8 процеса поради по-ниските комуникационни разходи.

    \item \textbf{Ефективността намалява} с броя процеси поради комуникационни разходи.
\end{enumerate}

\subsection{Влияние на грануларността}

При по-фина грануларност комуникационните разходи доминират, докато при много едра грануларност няма достатъчна гъвкавост за балансиране.

\section{Заключение}

Проектът демонстрира предимствата на P2P архитектурите:

\begin{enumerate}
    \item P2P елиминира тясното място на централизираните системи.
    \item Верижната топология е подходяща за големи системи.
    \item Пълният граф предлага по-бързо балансиране за малки клъстери.
    \item Изборът на топология зависи от конкретния сценарий.
\end{enumerate}

\section*{Източници}

\begin{enumerate}
    \item Wilkinson, B., \& Allen, M. (2004). \textit{Parallel Programming}. Prentice Hall.
    \item Gropp, W., et al. (2014). \textit{Using MPI}. MIT Press.
    \item Kumar, V., et al. (2003). \textit{Introduction to Parallel Computing}. Addison-Wesley.
    \item Amazon Web Services. (2024). \textit{EC2 Placement Groups}. AWS Documentation.
    \item Blumofe, R. D., \& Leiserson, C. E. (1999). Scheduling by work stealing. \textit{JACM}, 46(5).
\end{enumerate}

\end{document}
